\subsection{Complete design: a furrow system for 12 ha of silty clay loam}
\exbox{Design of a furrow system for 12 ha}{A farm of 12 ha of silty clay loam S10 is to be furrow irrigated. Furrows are 250 m long at 1.5 m spacing on a slope of 0.25 per cent, the wetted width is $w=0.75$ m, Manning $n=0.04$, and the net application is 75 mm at a 50 per cent trigger. A pump delivers 30 L\,s$^{-1}$. Derive the intake from the texture, compare four furrow inflows, choose one, examine a cutback, and find the number of furrows, the sets and the irrigation time.}\sol \givenline{$w=0.75$ m ; $n=0.04$}
\ans Soil S10 (sand 10, clay 35 per cent, organic matter 1.5 per cent): Saxton and Rawls give $\theta_{1500}=0.2130$, $\theta_{33}=0.3785$, $\theta_s=0.4903$ and $K_s=3.8$ mm\,h$^{-1}$
\ans Wetting-front suction $\psi_f=85.5$ mm from $\alpha=0.010$ cm$^{-1}$, $n=1.23$; moisture deficit at the 50 per cent trigger $\Delta\theta=0.1700$
\ans Green--Ampt cumulative intake at 10, 60 and 600 min: 4.7, 13.3 and 63.2 mm
\ans Log--log fit between 10 and 600 min: $Z=kt^{A}$ with $k=0.00103$ m\,min$^{-A}$, $A=0.634$ ($R^2=0.9980$)
\ans Net depth $Z_n=75$ mm, so the required contact time is \[\tau_n=(Z_n/k)^{1/A}=(0.075/0.00103)^{1/0.634}=871\] min (14.5 h)
\ans Trial $q_0=0.8$ L\,s$^{-1}$: Manning depth $y=8.2$ cm, inlet area $A_0=68$ cm$^2$; advance $t_L=68$ min, cut-off $t_{co}=939$ min; $E_a=31.2$, $DP=0.9$, $RO=67.9$ per cent, $CU=98.8$ per cent
\ans Trial $q_0=1.2$ L\,s$^{-1}$: Manning depth $y=9.6$ cm, inlet area $A_0=92$ cm$^2$; advance $t_L=45$ min, cut-off $t_{co}=916$ min; $E_a=21.3$, $DP=0.4$, $RO=78.3$ per cent, $CU=99.2$ per cent
\ans Trial $q_0=1.6$ L\,s$^{-1}$: Manning depth $y=10.7$ cm, inlet area $A_0=114$ cm$^2$; advance $t_L=36$ min, cut-off $t_{co}=907$ min; $E_a=16.1$, $DP=0.2$, $RO=83.6$ per cent, $CU=99.3$ per cent
\ans Trial $q_0=2$ L\,s$^{-1}$: Manning depth $y=11.6$ cm, inlet area $A_0=135$ cm$^2$; advance $t_L=32$ min, cut-off $t_{co}=902$ min; $E_a=13.0$, $DP=0.2$, $RO=86.9$ per cent, $CU=99.4$ per cent
\ans Choice: the highest $E_a$ with the requirement met ($E_s\ge0.99$) is $q_0=0.8$ L\,s$^{-1}$ per furrow, $E_a=31.2$ per cent
\ans Cutback to half the inflow after the front reaches the end: $V_{app}=24.17$ m$^3$ per furrow, $E_a=58.2$ per cent, $RO=40.1$ per cent, $E_s=1.000$ (an improvement of 27.0 points)
\ans Furrows per hectare at 1.5 m spacing and $L=250$ m: $10\,000/(1.5\times250)=26.7$; for 12 ha, $320$ furrows
\ans Set the water on 37 furrows at a time so that the supply $Q=29.6$ L\,s$^{-1}$ per set is within a 30 L\,s$^{-1}$ pump: $37\times0.8=29.6$ L\,s$^{-1}$
\ans Number of sets $=320/37=8.6$; each set runs $t_{co}=15.6$ h, so the whole field takes $135$ h ($6.8$ days of 20 h)
\ans Gross volume for the field $=320\times45.07=14422$ m$^3$, or 120 mm over 12 ha \finans


